% Figure: pull force of a solenoid actuator against air-gap length. The
% force rises steeply as the gap closes because the gap reluctance, and
% hence the energy gradient, scales as 1/g^2 at fixed magnetomotive force.
\begin{figure}[htbp]
\centering
\begin{tikzpicture}[>=Latex]
\begin{axis}[
  width=11.5cm, height=7cm,
  xlabel={air gap $g$ ($\si{\milli\meter}$)},
  ylabel={pull force $F$ ($\si{\newton}$)},
  xmin=0, xmax=10,
  ymin=0, ymax=260,
  xtick={0,2,4,6,8,10},
  legend pos=north east,
  legend cell align=left,
  font=\scriptsize,
  grid=both,
  grid style={enggray!20},
]
% F = (mu0 (NI)^2 A) / (2 g^2); clamp near g->0 by starting at g=0.5 mm.
% Use constant K so that F(1mm)=50 N => K = 50. Then F = 50/g^2 (g in mm).
\addplot[engblue, line width=1.2pt, domain=0.5:10, samples=120]
  {50/(x*x)};
\addlegendentry{$F = \dfrac{\mu_0 (NI)^2 A}{2\,g^2}$}
% mark the working stroke endpoints
\addplot[engred, only marks, mark=*, mark size=2pt]
  coordinates {(8,0.78) (1,50)};
\draw[enggray, dashed] (axis cs:1,0) -- (axis cs:1,50);
\draw[enggray, dashed] (axis cs:8,0) -- (axis cs:8,0.78);
\node[engred, anchor=south west, font=\scriptsize] at (axis cs:1.1,52)
  {seated, $g = 1\,\si{\milli\meter}$};
\node[engred, anchor=west, font=\scriptsize] at (axis cs:8.1,12)
  {open, $g = 8\,\si{\milli\meter}$};
\end{axis}
\end{tikzpicture}
\caption[Solenoid-actuator pull force against air gap]{Pull force of a
solenoid actuator against the length of the air gap between the plunger
and the seat. At fixed magnetomotive force the force scales as
$1/g^2$, so it is weakest when the plunger is far out and the work must
begin, and strongest as the plunger seats. A design that must pull a
load over a long stroke therefore fights the worst force at the worst
moment, which is why long-stroke actuators use a stepped or conical pole
face to flatten the curve, and why a return spring sized for the seated
force will not be overcome when the gap is open. The two marked points
are the open and seated ends of an $8\,\si{\milli\meter}$ working
stroke; the force across it varies by a factor of sixty.}
\label{fig:vol04ch08:actuator-force}
\end{figure}
